\begin{abstractAr}
في سياق الثورة الصناعية الرابعة وانتشار أنظمة التصنيع الإضافي، تُستخدم الطابعات ثلاثية الأبعاد بتقنية النمذجة بالترسيب المنصهر (FDM) بشكل متزايد في بيئات الإنتاج، والنمذجة السريعة، والبحث العلمي. إلا أن موثوقيتها التشغيلية لا تزال محدودة بسبب الأعطال المتكررة مثل انسداد الفوهة، والاهتزازات غير الطبيعية لمحاور الحركة، وظاهرة الارتفاع المفرط في الحرارة، مما يؤدي إلى توقف غير مخطط له وتدهور جودة الطباعة.
يقدم هذا العمل تصميم وتنفيذ توأم رقمي هجين للصيانة التنبؤية لطابعة ثلاثية الأبعاد من نوع BCN3D+. الهدف الرئيسي هو تمكين المراقبة في الوقت الفعلي والتنبؤ المبكر بالأعطال من أجل تقليل انقطاعات الإنتاج وتحسين عمر الماكينة.
يعتمد النظام المقترح على بنية متعددة المستويات تجمع بين النمذجة الهندسية والسلوكية، ويتم مزامنتها في الوقت الفعلي عبر بنية تحتية تعتمد على إنترنت الأشياء (IoT). تشتمل التجهيزات المادية على مستشعر اهتزاز MPU-6050 ومستشعر حرارة KY-013، متصلين بمتحكم دقيق ESP8266 مسؤول عن جمع البيانات وإرسالها. بالإضافة إلى ذلك، يُستخدم فأرة بصرية معدَّلة الهيكل والمشغِّل كمشفِّر لخيط الطباعة.
تتبع معالجة البيانات نهجاً هرمياً هجيناً: قواعد منطقية حتمية (المستوى 0)، وكشف الحالات الشاذة باستخدام خوارزمية Isolation Forest (المستوى 1)، وتصنيف خاضع للإشراف باستخدام XGBoost (المستوى 2). يتيح هذا النهج المتدرج استراتيجية كشف تدريجية تتراوح من تحديد الحالات الشاذة البسيطة إلى التنبؤ المتقدم بالأعطال.
تم التحقق من صحة النظام من خلال ثلاث دراسات حالة تجريبية: انسداد الفوهة، وشذوذ اهتزازات محاور الحركة، وظروف الارتفاع المفرط في الحرارة. تظهر النتائج دقة إجمالية مرجحة بنسبة 90.8%، وانخفاضاً في معدل فشل الطباعة بنسبة 67%، وتقليصاً في وقت التوقف غير المخطط له بنسبة 46%. تظل التكلفة الإجمالية للأجهزة أقل من 50 يورو، مما يجعل الحل في متناول الشركات الصغيرة والمتوسطة ومساحات التصنيع المشتركة (الميكرسبيس).
\end{abstractAr}

\begin{keywordsAr}
توأم رقمي، صيانة تنبؤية، طباعة ثلاثية الأبعاد، نمذجة بالترسيب المنصهر (FDM)، إنترنت الأشياء، تعلم آلي، الثورة الصناعية الرابعة.
\end{keywordsAr}